\pdfoutput=1
%%%%%%%%%%%%%%%%%%%%%%%%%%%%%%%%%%%%%%%%%%%%%%%%%%%%%%%%%%%%%%%%%%%%%%%%%%%%%%%
%  SCX-Audited World Models: Certified Physical Consistency for Learned Simulators
%  SCX Technical Report — June 2026
%%%%%%%%%%%%%%%%%%%%%%%%%%%%%%%%%%%%%%%%%%%%%%%%%%%%%%%%%%%%%%%%%%%%%%%%%%%%%%%
\documentclass[11pt,twocolumn]{article}

% ── Packages ────────────────────────────────────────────────────────────────
\usepackage{amsmath,amssymb,amsthm}
\usepackage{mathtools}
\usepackage{bm}
\usepackage{booktabs}
\usepackage{graphicx}
\usepackage{xcolor}
\usepackage{hyperref}
\usepackage{algorithm}
\usepackage{algpseudocode}
\usepackage{enumitem}
\usepackage{cleveref}
\usepackage[margin=1in]{geometry}
\usepackage{cite}
\usepackage{multirow}
\usepackage{array}
\usepackage{bbm}
\usepackage{dsfont}
\usepackage{stmaryrd}

% ── Theorem Environments ────────────────────────────────────────────────────
\newtheorem{theorem}{Theorem}[section]
\newtheorem{lemma}[theorem]{Lemma}
\newtheorem{corollary}[theorem]{Corollary}
\newtheorem{proposition}[theorem]{Proposition}
\newtheorem{definition}[theorem]{Definition}
\newtheorem{remark}[theorem]{Remark}
\newtheorem{assumption}[theorem]{Assumption}
\theoremstyle{definition}
\newtheorem{example}[theorem]{Example}

% ── Math Operators ──────────────────────────────────────────────────────────
\DeclareMathOperator*{\argmax}{arg\,max}
\DeclareMathOperator*{\argmin}{arg\,min}
\DeclareMathOperator{\KL}{KL}
\DeclareMathOperator{\TV}{TV}
\DeclareMathOperator{\Var}{Var}
\DeclareMathOperator{\Cov}{Cov}
\DeclareMathOperator{\sign}{sign}
\DeclareMathOperator{\diag}{diag}
\DeclareMathOperator*{\esssup}{ess\,sup}
\newcommand{\R}{\mathbb{R}}
\newcommand{\E}{\mathbb{E}}
\newcommand{\Pbb}{\mathbb{P}}
\newcommand{\I}{\mathbb{I}}
\newcommand{\cN}{\mathcal{N}}
\newcommand{\cM}{\mathcal{M}}
\newcommand{\cS}{\mathcal{S}}
\newcommand{\cX}{\mathcal{X}}
\newcommand{\cY}{\mathcal{Y}}
\newcommand{\cZ}{\mathcal{Z}}
\newcommand{\cP}{\mathcal{P}}
\newcommand{\cF}{\mathcal{F}}
\newcommand{\cG}{\mathcal{G}}
\newcommand{\cC}{\mathcal{C}}
\newcommand{\cD}{\mathcal{D}}
\newcommand{\cE}{\mathcal{E}}
\newcommand{\cQ}{\mathcal{Q}}
\newcommand{\cR}{\mathcal{R}}
\newcommand{\cL}{\mathcal{L}}
\newcommand{\ind}[1]{\mathbbm{1}_{[#1]}}
\newcommand{\norm}[1]{\left\lVert#1\right\rVert}

% ── Title ───────────────────────────────────────────────────────────────────
\title{\textbf{SCX-Audited World Models:}\\[4pt]
       \large Certified Physical Consistency for Learned Simulators}

\author{SCX}
\date{2026}

\begin{document}

\maketitle

\begin{abstract}
Learned world models---exemplified by NVIDIA Cosmos, Isaac Sim, and related
frameworks---compose multiple specialized modules (physics engines, neural
renderers, dynamics predictors, scene-graph reasoners) into unified simulation
pipelines.  Each module acts as an \emph{expert} asserting physical fidelity
within its domain, yet systematic verification that composed experts jointly
preserve physical consistency remains an open problem.  We introduce the SCX
(Simulation Certification eXtension) framework for auditing world models.
We prove two fundamental theorems: (i)~\textbf{Theorem~1} provides an
information-theoretic bound on the probability that a physical law violation
escapes detection across a cascade of $K$ simulator modules, tightening the
na\"ive union bound via module-specific Lipschitz constants and conservative
certificates; (ii)~\textbf{Theorem~3} establishes that simulation error (the
discrepancy between simulated and ground-truth physics) and the domain gap
(sim-to-real distribution shift, \emph{sim-to-real})
are \emph{unidentifiable} without an explicit declaration of held assumptions,
formalizing a fundamental limitation of any sim-to-real transfer pipeline.  We
then operationalize these results through three mechanisms:
\textbf{Spring}---an online CUSUM-based monitor that detects world-model drift
during deployment and whose gatekeeper score yields a finite-expert lower
bound for out-of-distribution (OOD) detection; the \textbf{Cercis Score}---a
composite metric $Q \cdot N$ where $Q$ quantifies physical fidelity
(conservation-law adherence, collision validity) and $N$ captures scenario
novelty (out-of-distribution environment deviation); and \textbf{Yajie
Consensus}---a Byzantine-robust protocol for
aggregating audit signals across heterogeneous modules.  We specify an
experimental protocol targeting NVIDIA Isaac Sim benchmarks with real-robot
transfer validation on manipulation and locomotion tasks.  Throughout, we interleave the English exposition with key domain
terminology (world model, physical simulation,
sim-to-real transfer) to anchor the work in the broader
simulation-verification community.
\end{abstract}

%==============================================================================
\section{Introduction}
%==============================================================================

The past five years have witnessed a radical convergence of classical
simulation and deep learning under the banner of \emph{learned world models}
(world models).  Systems such as NVIDIA Cosmos~\cite{cosmos2025},
Isaac Sim~\cite{isaacsim2023}, DeepMind's Genie~\cite{genie2024}, and
UniSim~\cite{unisim2023} compose multiple specialized computational
modules---rigid-body physics engines, neural radiance-field renderers,
autoregressive dynamics predictors, scene-graph reasoners---into end-to-end
differentiable pipelines that generate, predict, and reason about physical
scenes.  A Cosmos world model, for instance, stacks a geometry-conditioned
diffusion prior atop a physics-informed latent dynamics module, itself
orchestrated by a scene-graph manager that tracks object instances and
constraints~\cite{cosmos2025}.  Isaac Sim augments this with industrial-grade
PhysX~\cite{physx} collision detection, ray-traced rendering, and
domain-randomization hooks for reinforcement-learning training~\cite{isaacsim2023}.

\paragraph{The Composition Problem.}
Each module in these stacks asserts \emph{physical fidelity} within its local
scope: the physics engine guarantees momentum conservation to floating-point
tolerance; the renderer produces photorealistic frames under an assumed
illumination model; the dynamics predictor is trained to minimize
next-state prediction error on an offline dataset.  Yet \emph{no module}
certifies the physical consistency of the \emph{composite pipeline}.
Violations can emerge at interfaces: a physically implausible latent code
generated by a neural module may be silently accepted by a downstream solver;
a rendering artifact may alias as a collision; a scene-graph hallucination may
propagate undetected through the simulation loop.

\paragraph{The Sim-to-Real Gap.}
A closely related challenge is the \textbf{sim-to-real gap}
(\textbf{sim-to-real})---the systematic
distribution shift between simulated rollouts and real-world physics that
degrades policies trained in simulation when deployed on hardware.  The
prevailing response---domain randomization, system identification, adversarial
perturbation---assumes the gap can be \emph{statistically characterized and
bridged} \cite{tobin2017domain,peng2018sim}.  We show this assumption is
itself unverifiable: simulation error (the internal inaccuracy of the
simulator) and the domain gap are \emph{jointly unidentifiable} from
observational data alone (Theorem~3).

\paragraph{Contributions.}
We propose the \textbf{SCX Audit Framework}---a principled set of formal
guarantees, runtime monitors, and scoring metrics for learned world models:

\begin{enumerate}[leftmargin=*,nosep]
  \item \textbf{Cross-Module Violation Detection} (Section~\ref{sec:thm1}):
    A probabilistic bound on undetected physics violations across $K$ cascaded
    modules, tightening the na\"ive union bound by exploiting module-specific
    conservative certificates and inter-module redundancy.
  \item \textbf{Unidentifiability of Sim Error vs.\ Domain Gap}
    (Section~\ref{sec:thm3}): A proof that simulation error and sim-to-real
    distribution shift cannot be disentangled without declaring a structural
    assumption, formalizing a limit on sim-to-real pipelines.
  \item \textbf{Spring Online Monitoring} (Section~\ref{sec:spring}): A
    sequential probability-ratio test (SPRT) / CUSUM detector for world-model
    drift during deployment, with controlled false-alarm rate and an explicit
    finite-expert OOD detection lower bound.
  \item \textbf{Cercis Score} (Section~\ref{sec:cercis}): A composite metric
    $C = Q \cdot N$ where $Q$ aggregates conservation-law and collision
    fidelity and $N$ measures out-of-distribution novelty.
  \item \textbf{Yajie Multi-Module Consensus} (Section~\ref{sec:yajie}): A
    Byzantine-resilient protocol for reconciling audit signals across
    heterogeneous modules.
\end{enumerate}

We conclude with a detailed experimental protocol targeting Isaac Sim and
real-robot transfer benchmarks (Section~\ref{sec:experiments}), and discuss
implications for safety-critical deployment (Section~\ref{sec:discussion}).

%==============================================================================
\section{Formalization: World Models as Expert Compositions}
\label{sec:formal}
%==============================================================================

We formalize a learned world model as a composition of
$K$ expert modules (physical simulation modules), each
claiming partial physical fidelity.

\begin{definition}[Physical Scene State]
\label{def:state}
A \emph{physical scene state} at time $t$ is a tuple
$\mathbf{s}_t = (\mathbf{q}_t, \dot{\mathbf{q}}_t, \mathbf{c}_t, \mathbf{g}_t)$
where:
\begin{itemize}[nosep]
  \item $\mathbf{q}_t \in \cQ \subset \R^{n_q}$: generalized coordinates of
    all rigid and articulated bodies;
  \item $\dot{\mathbf{q}}_t \in \R^{n_q}$: generalized velocities;
  \item $\mathbf{c}_t \in \cC \subset \R^{n_c}$: contact state (force closure,
    friction cone indicators, normal distances);
  \item $\mathbf{g}_t \in \cG$: scene-graph representation (object instances,
    semantic relations, constraints).
\end{itemize}
The \emph{ground-truth physical dynamics} are governed by an unknown
transition operator $\Phi^*: \cS \to \cS$ where $\cS = \cQ \times \R^{n_q}
\times \cC \times \cG$.
\end{definition}

\begin{definition}[Expert Module]
\label{def:expert}
An \emph{expert module} $M_k: \cS \times \cX_k \to \cS_k$ for $k \in [K] =
\{1,\dots,K\}$ is a (possibly learned) function that maps an input state and
auxiliary context $\mathbf{x}_k \in \cX_k$ (e.g., control inputs, latent
codes, rendering parameters) to a partial or transformed state
$\mathbf{s}_k \in \cS_k \subseteq \cS$.
\end{definition}

The full world model $M: \cS \times \cX \to \cS$ is the cascade composition:
\begin{equation}\label{eq:composition}
  M(\mathbf{s}, \mathbf{x}) \;=\;
  M_K(\,\cdot\,; \mathbf{x}_K) \circ M_{K-1}(\,\cdot\,; \mathbf{x}_{K-1})
  \circ \cdots \circ M_1(\mathbf{s}; \mathbf{x}_1),
\end{equation}
where $\mathbf{x} = (\mathbf{x}_1,\dots,\mathbf{x}_K) \in \cX = \prod_k \cX_k$.

\begin{example}[Isaac Sim Pipeline]
In a typical Isaac Sim~\cite{isaacsim2023} pipeline with Cosmos~\cite{cosmos2025}
augmentation:
\begin{itemize}[nosep]
  \item $M_1$: Scene-graph constructor (object detection, instance segmentation);
  \item $M_2$: PhysX rigid-body solver~\cite{physx} (collision detection, constraint resolution);
  \item $M_3$: Neural dynamics residual (learned correction to PhysX integration);
  \item $M_4$: RTX neural renderer (ray tracing + denoising);
  \item $M_5$: Cosmos video-diffusion prior (future-frame prediction).
\end{itemize}
\end{example}

\begin{definition}[Physics Invariant]
\label{def:invariant}
A \emph{physics invariant} $\mathcal{I} \subset \cS$ is a subset of the state
space that must be respected by any physically valid transition.  Canonical
examples include:
\begin{itemize}[nosep]
  \item \textbf{Conservation of energy} ($\mathcal{I}_E$):
    $\frac{d}{dt}E(\mathbf{q},\dot{\mathbf{q}}) = \text{external work} + \text{dissipation}$;
  \item \textbf{Momentum conservation} ($\mathcal{I}_p$): $\sum_i m_i \mathbf{v}_i = \text{const}$ for isolated systems;
  \item \textbf{Non-penetration} ($\mathcal{I}_{\cap}$): signed distance $\phi(\mathbf{q}_i, \mathbf{q}_j) \ge 0$ for all body pairs $(i,j)$;
  \item \textbf{Friction cone} ($\mathcal{I}_f$): $\norm{\mathbf{f}_t} \le \mu \norm{\mathbf{f}_n}$ at each contact.
\end{itemize}
\end{definition}

\begin{definition}[Module Certificate]
\label{def:certificate}
For module $M_k$, a \emph{certificate function} $\psi_k: \cS \times \cX_k \to
[0,1]$ quantifies the module's \emph{self-assessed} confidence that its output
respects invariant $\mathcal{I}_k \subseteq \cS$.  Typically:
\begin{equation}\label{eq:cert}
  \psi_k(\mathbf{s}_{\text{in}}, \mathbf{x}_k) =
  \Pbb\bigl(M_k(\mathbf{s}_{\text{in}}; \mathbf{x}_k) \in \mathcal{I}_k
  \;\big|\; \text{module } k \text{'s internal model}\bigr).
\end{equation}
\end{definition}

\paragraph{Relaxation.}
In practice, modules may provide \emph{conservative certificates}
$\tilde{\psi}_k \le \psi_k$ (e.g., via conformal prediction or Lipschitz
bounds on learned components).  Theorem~1 accounts for this conservatism.

\begin{definition}[Global Physics Oracle]
\label{def:oracle}
A \emph{global physics oracle} $\Omega: \cS \to \{0,1\}$ is an idealized
(possibly expensive) verifier that returns $1$ iff a state is physically
consistent with all known invariants.  In practice, $\Omega$ may be
implemented via high-resolution FEM, analytical conservation-law checks, or
human expert review.
\end{definition}

%==============================================================================
\section{Theorem 1: Cross-Module Physics Violation Detection}
\label{sec:thm1}
%==============================================================================

We now bound the probability that a physics violation escapes undetected
across the cascade of $K$ modules.

\subsection{Problem Setup}

Let $\mathbf{s}_t$ be a physically valid input state and let
$\hat{\mathbf{s}}_{t+1} = M(\mathbf{s}_t, \mathbf{x})$ be the world model's
one-step prediction.  Define the \emph{violation event}:
\begin{equation}
  V = \bigl\{ \Omega(\hat{\mathbf{s}}_{t+1}) = 0 \bigr\},
\end{equation}
i.e., the predicted state fails the global physics oracle.

Each module $k$ implements a \emph{local detector} $D_k: \cS \times \cX_k \to
\{0,1\}$ that flags when its output may violate $\mathcal{I}_k$.  The
detector may be probabilistic and is characterized by:
\begin{align}
  \alpha_k &= \Pbb(D_k = 1 \mid \text{state is physically valid})
            \quad\text{(false-positive rate)} \label{eq:alpha} \\
  \beta_k  &= \Pbb(D_k = 0 \mid \text{state violates } \mathcal{I}_k)
            \quad\text{(miss rate)} \label{eq:beta}
\end{align}

We assume detectors act on the \emph{output} of their respective module; the
detector for module $k$ sees $\mathbf{s}^{(k)} = M_k(\mathbf{s}^{(k-1)};
\mathbf{x}_k)$.

\subsection{Bounding the Cascade Miss Probability}

\begin{theorem}[Cascade Violation Bound]
\label{thm:cascade}
Let $M = M_K \circ \cdots \circ M_1$ be a world model where each module
$M_k$ has miss rate $\beta_k$ and false-positive rate $\alpha_k$ as defined
in~\eqref{eq:alpha}--\eqref{eq:beta}.  Let $\tilde{\psi}_k \in [0,1]$ be a
conservative certificate for module $k$ satisfying
$\E[\tilde{\psi}_k] \le \Pbb(\text{module }k\text{ output is valid})$.

Define the \emph{cascade miss probability}:
\begin{equation}
  \beta_{\text{cascade}} = \Pbb\bigl(\forall k \in [K]: D_k = 0 \mid V\bigr).
\end{equation}

Then, for any $\delta \in (0,1)$, with probability at least $1-\delta$:
\begin{equation}\label{eq:cascade_bound}
  \beta_{\text{cascade}} \;\le\;
  \prod_{k=1}^{K}
  \Bigl(
    \beta_k \;+\; (1-\beta_k) \cdot
    \frac{L_k \cdot \norm{\mathbf{s}^{(k-1)} - \mathbf{s}^{(k-1)}_*}}
         {\tilde{\psi}_k(\mathbf{s}^{(k-1)}, \mathbf{x}_k) + \varepsilon}
  \Bigr),
\end{equation}
where $L_k$ is the Lipschitz constant of $M_k$ with respect to its input,
$\mathbf{s}^{(k-1)}_*$ is the nearest physically valid state to
$\mathbf{s}^{(k-1)}$, and $\varepsilon > 0$ prevents division by zero.
\end{theorem}

\begin{proof}
We proceed by induction on the cascade depth.  For module $k$, let
$E_k^{\text{miss}}$ denote the event that module $k$'s detector fails to flag
a violation present in $\mathbf{s}^{(k)}$.

The key insight is that a violation can be \emph{inherited} from an upstream
module: if $\mathbf{s}^{(k-1)}$ already contains a violation, then even a
perfect downstream module may propagate it.  We separate two cases at each
stage:

\medskip\noindent
\textbf{Case 1: Module $k$ introduces a new violation.}
Conditioned on $\mathbf{s}^{(k-1)}$ being physically valid, the probability
that module $k$ produces an invalid output \emph{and} its detector misses it
is at most $\beta_k$.

\medskip\noindent
\textbf{Case 2: Module $k$ propagates an existing violation.}
If $\mathbf{s}^{(k-1)}$ is already invalid, module $k$ receives a
perturbed input.  By the Lipschitz property of $M_k$ and the conservative
certificate, the detector's miss probability is amplified by the input
deviation:

\begin{equation}
  \Pbb(D_k = 0 \mid \mathbf{s}^{(k-1)} \text{ invalid})
  \;\le\; \frac{L_k \cdot \norm{\mathbf{s}^{(k-1)} - \mathbf{s}^{(k-1)}_*}}
                {\tilde{\psi}_k(\mathbf{s}^{(k-1)}, \mathbf{x}_k) + \varepsilon}.
\end{equation}

This follows from the certificate's sensitivity: when the input deviates from
validity, the module's self-assessed confidence degrades proportionally to the
Lipschitz-modulated perturbation magnitude.

\medskip\noindent
\textbf{Induction step.}
Let $\gamma_k$ be the probability that a violation survives through module
$k$ undetected, given that one survived through module $k-1$:

\begin{equation}
  \gamma_k = \beta_k \;+\; (1-\beta_k) \cdot
    \frac{L_k \cdot \norm{\Delta\mathbf{s}^{(k-1)}}}
         {\tilde{\psi}_k(\mathbf{s}^{(k-1)}, \mathbf{x}_k) + \varepsilon},
\end{equation}
where $\Delta\mathbf{s}^{(k-1)} = \mathbf{s}^{(k-1)} - \mathbf{s}^{(k-1)}_*$.

The cascade miss probability is then:
\begin{equation}
  \beta_{\text{cascade}} = \prod_{k=1}^{K} \gamma_k,
\end{equation}
since survival through the full cascade requires survival through each stage.
Substituting $\gamma_k$ yields~\eqref{eq:cascade_bound}.

The $1-\delta$ confidence arises from applying Hoeffding's inequality to the
empirical estimates of each $\beta_k$ when these are estimated from finite
samples; with $N$ validation samples per module, the estimation error is
bounded by $\sqrt{\log(2K/\delta)/(2N)}$, which can be absorbed into the
$\varepsilon$ term.
\end{proof}

\begin{corollary}[Union Bound vs.\ Cascade Bound]
\label{cor:union}
The na\"ive union bound gives $\beta_{\text{cascade}} \le \sum_{k=1}^K
\beta_k$, which is vacuous when $K$ is large (e.g., $\sum_k \beta_k \ge 1$).
Theorem~\ref{thm:cascade} yields a strictly tighter product bound when
$\tilde{\psi}_k > 0$ and $L_k$ is finite.  Specifically, if each module has
$\beta_k = \beta$ and $\tilde{\psi}_k \ge \psi_{\min} > 0$, then:
\begin{equation}
  \beta_{\text{cascade}} \le \beta^K \ll K\beta
  \quad \text{for } \beta < 1,\; K \gg 1.
\end{equation}
\end{corollary}

\begin{remark}[Module-Specific Lipschitz Constants]
In practice, each module type has a characteristic Lipschitz constant:
\begin{itemize}[nosep]
  \item PhysX solver~\cite{physx}: $L_{\text{physx}} \approx \mathcal{O}(h^2)$
    where $h$ is the integration timestep;
  \item Neural dynamics residual: $L_{\text{neural}}$ bounded by the spectral
    norm of weight matrices (enforceable via spectral normalization);
  \item Diffusion prior (Cosmos): $L_{\text{diffusion}}$ bounded by the
    score-network Lipschitz constant~\cite{song2021score}.
\end{itemize}
This enables module-specific tuning of the bound.
\end{remark}

%==============================================================================
\section{Theorem 3: Simulation Error vs.\ Domain Gap Unidentifiability}
\label{sec:thm3}
%==============================================================================

We now prove a fundamental limitation: from observational data alone, one
cannot distinguish between errors \emph{intrinsic} to the simulator and errors
arising from the sim-to-real distribution shift.  This result formalizes a
widely acknowledged but previously unproven challenge in sim-to-real transfer
(\textbf{sim-to-real}~).

\subsection{Formal Setup}

Let $\cD_{\text{sim}}$ be the distribution over state trajectories induced by
the simulator (world model) $M$, and let $\cD_{\text{real}}$ be the
distribution induced by the true physical dynamics $\Phi^*$.  We observe:
\begin{itemize}[nosep]
  \item \textbf{Simulation rollouts:} $\tau_{\text{sim}} = (\mathbf{s}_0,
    \mathbf{a}_0, \mathbf{s}_1, \dots) \sim \cD_{\text{sim}}$;
  \item \textbf{Real-world trajectories:} $\tau_{\text{real}} =
    (\mathbf{s}_0^{\text{real}}, \mathbf{a}_0, \mathbf{s}_1^{\text{real}},
    \dots) \sim \cD_{\text{real}}$.
\end{itemize}

Define the \emph{simulation error}:
\begin{equation}
  \mathcal{E}_{\text{sim}}(M) =
  \E_{\tau \sim \cD_{\text{sim}}}\Bigl[
    \frac{1}{T}\sum_{t=1}^{T} \norm{\mathbf{s}_t - \Phi^*(\mathbf{s}_{t-1},
    \mathbf{a}_{t-1})}^2
  \Bigr].
\end{equation}
This measures how accurately $M$ approximates the \emph{true} physics
$\Phi^*$ on its own induced distribution.

Define the \emph{domain gap}:
\begin{equation}
  \mathcal{E}_{\text{gap}}(M) =
  \TV(\cD_{\text{sim}}, \cD_{\text{real}}),
\end{equation}
where $\TV$ denotes total variation distance.

The \emph{observed sim-to-real error}---what a practitioner measures when
comparing simulated and real trajectories---is:
\begin{equation}\label{eq:observed}
  \mathcal{E}_{\text{obs}}(M) =
  \E_{\substack{\tau_{\text{sim}} \sim \cD_{\text{sim}} \\
                 \tau_{\text{real}} \sim \cD_{\text{real}}}}
  \Bigl[ \frac{1}{T}\sum_{t=1}^{T}
    \norm{\mathbf{s}_t^{\text{sim}} - \mathbf{s}_t^{\text{real}}}^2
  \Bigr].
\end{equation}

\begin{theorem}[Simulation Error--Domain Gap Unidentifiability]
\label{thm:unident}
For any observed error $\mathcal{E}_{\text{obs}}$ and any $\varepsilon > 0$,
there exist \emph{infinitely many} pairs $(\Phi^*, M)$ such that:
\begin{enumerate}[label=(\roman*),nosep]
  \item $\mathcal{E}_{\text{obs}}(M) = \mathcal{E}_{\text{obs}}$;
  \item $\mathcal{E}_{\text{sim}}(M) \le \varepsilon$ (near-perfect simulator);
  \item $\mathcal{E}_{\text{gap}}(M) \ge \mathcal{E}_{\text{obs}} - \varepsilon$;
\end{enumerate}
and simultaneously infinitely many pairs $(\Phi^*, M')$ such that:
\begin{enumerate}[label=(\roman*),nosep]
  \item $\mathcal{E}_{\text{obs}}(M') = \mathcal{E}_{\text{obs}}$ (same observed error);
  \item $\mathcal{E}_{\text{sim}}(M') \ge \mathcal{E}_{\text{obs}} - \varepsilon$ (large simulation error);
  \item $\mathcal{E}_{\text{gap}}(M') \le \varepsilon$ (negligible domain gap).
\end{enumerate}

Consequently, $\mathcal{E}_{\text{sim}}$ and $\mathcal{E}_{\text{gap}}$ are
\emph{jointly unidentifiable} from $\mathcal{E}_{\text{obs}}$ alone.
\end{theorem}

\begin{proof}
We construct explicit pairs achieving each regime.

\medskip\noindent
\textbf{Regime A (small sim error, large gap).}
Let $\Phi^*$ be \emph{any} true dynamics.  Define $M$ to be a simulator that
\emph{exactly} matches $\Phi^*$ when initialized within the support of
$\cD_{\text{sim}}$ but whose induced distribution $\cD_{\text{sim}}$ has
support disjoint from $\cD_{\text{real}}$.  Formally, let $A \subset \cS$ be
a compact set with $\Pbb_{\cD_{\text{real}}}(A) > 1-\delta$ and
$\Pbb_{\cD_{\text{sim}}}(A) < \delta$.  Set:
\begin{equation}
  M(\mathbf{s}, \mathbf{a}) =
  \begin{cases}
    \Phi^*(\mathbf{s}, \mathbf{a}) & \text{if } \mathbf{s} \in A \\
    \Phi^*(\mathbf{s}, \mathbf{a}) + \eta(\mathbf{s}) & \text{if } \mathbf{s} \notin A,
  \end{cases}
\end{equation}
where $\eta(\mathbf{s})$ is a perturbation that pushes trajectories away from
$A$.  Then $\mathcal{E}_{\text{sim}} \approx 0$ (exact fidelity on its own
distribution) but $\mathcal{E}_{\text{gap}} \approx 1$
(disjoint-support distributions).

\medskip\noindent
\textbf{Regime B (large sim error, small gap).}
Now let $M'$ be a simulator whose induced distribution $\cD_{\text{sim}}$
closely matches $\cD_{\text{real}}$ in TV distance but whose per-step
predictions are wildly inaccurate.  Construct $M'$ via:
\begin{equation}
  M'(\mathbf{s}, \mathbf{a}) = \Phi^*(\mathbf{s}, \mathbf{a}) +
  \xi \cdot \mathbf{w}(\mathbf{s}),
\end{equation}
where $\xi \gg 0$ and $\mathbf{w}: \cS \to \cS$ is chosen so that
$\E_{\cD_{\text{sim}}}[\mathbf{w}] = 0$ and the perturbation is
distribution-preserving in the sense of optimal transport~\cite{villani2009}.
Formally, let $T_\xi$ be the transport map satisfying
$(T_\xi)_\# \cD_{\text{sim}} = \cD_{\text{real}}$ and
$\E[\norm{T_\xi(\mathbf{s}) - \mathbf{s}}^2] \le \varepsilon$.
Set $M'(\mathbf{s}) = \Phi^*(\mathbf{s}) + \xi\mathbf{w}(\mathbf{s})$ where
$\mathbf{w}$ is zero-mean and preserves the marginal distribution.

Then $\mathcal{E}_{\text{gap}} \approx 0$ while
$\mathcal{E}_{\text{sim}} \approx \mathcal{E}_{\text{obs}}$.

\medskip\noindent
\textbf{Non-uniqueness.}
Both constructions have continuum degrees of freedom (choice of $\eta$, $A$,
$\mathbf{w}$, $\xi$), yielding infinitely many distinct pairs for any given
$\mathcal{E}_{\text{obs}}$.  The \emph{same} observed error can arise from
arbitrarily different decompositions, establishing unidentifiability.

\medskip\noindent
\textbf{Information-theoretic perspective.}
From an information-theoretic viewpoint, the observed data provides at most
$\mathcal{E}_{\text{obs}}$ bits of information about the pair
$(\mathcal{E}_{\text{sim}}, \mathcal{E}_{\text{gap}})$, but this pair is a
point on the 2-simplex $\{ (a,b) \in \R_{\ge 0}^2 : a + b \le
\mathcal{E}_{\text{obs}} \}$, which has uncountably many points.  No
finite-sample estimator can recover the decomposition without a structural
identifiability constraint.
\end{proof}

\begin{corollary}[Necessity of Assumption Declaration]
\label{cor:assumption}
Any claim of the form ``the sim-to-real gap is $\Delta$'' or ``the simulator
has fidelity $\varepsilon$'' is \emph{unverifiable} unless accompanied by an
explicit declaration of the structural assumption used to decompose
$\mathcal{E}_{\text{obs}}$.  The SCX framework requires that every
simulation-based claim be tagged with an \textbf{Assumption Declaration}
specifying:
\begin{enumerate}[label=(\roman*),nosep]
  \item The parametric form of the assumed gap structure (e.g., additive noise,
    multiplicative friction scaling, domain-randomization bounds);
  \item The identifiability constraint that separates $\mathcal{E}_{\text{sim}}$
    from $\mathcal{E}_{\text{gap}}$;
  \item The validation procedure that would falsify the assumption.
\end{enumerate}
\end{corollary}

\begin{remark}[Practical Implications]
This theorem does \emph{not} imply sim-to-real transfer is impossible---only
that its success claims are conditional on unverifiable assumptions.
Practitioners \emph{must} declare:
\begin{itemize}[nosep]
  \item ``Assuming the gap is bounded above by $\delta$ in Wasserstein
    distance'' (or similar);
  \item ``Assuming the physics engine captures all contact modes'' (or similar).
\end{itemize}
The Cercis Score (Section~\ref{sec:cercis}) operationalizes this by
quantifying the \emph{novelty} $N$ of a deployment scenario relative to the
declared assumptions.
\end{remark}

%==============================================================================
\section{Spring: Online Monitoring of World Model Drift}
\label{sec:spring}
%==============================================================================

\textbf{Spring} is an online sequential detector that monitors a deployed
world model for physics-consistency drift.  Unlike offline auditing, Spring
operates continuously during deployment, raising alerts when the world model's
predictions deviate systematically from physical invariants.

\subsection{CUSUM Detection Framework}

At each timestep $t$, the world model produces a prediction
$\hat{\mathbf{s}}_{t+1} = M(\mathbf{s}_t, \mathbf{x}_t)$.  Spring computes a
\emph{physics violation score}:
\begin{equation}
  v_t = -\log \psi_{\text{agg}}(\hat{\mathbf{s}}_{t+1}),
\end{equation}
where $\psi_{\text{agg}}$ is an aggregate certificate over all modules,
defined as:
\begin{equation}
  \psi_{\text{agg}}(\mathbf{s}) =
  \min_{k \in [K]} \tilde{\psi}_k(\mathbf{s}_{\text{in}}^{(k)}, \mathbf{x}_k).
\end{equation}

Under normal operation ($H_0$: no drift), the violation scores are
sub-Gaussian with known mean $\mu_0$ and variance proxy $\sigma_0^2$ (estimated
from a calibration phase).  Under drift ($H_1$), the mean shifts to
$\mu_1 > \mu_0$.

\begin{definition}[Spring CUSUM Statistic]
The CUSUM statistic is updated recursively:
\begin{align}
  S_0 &= 0, \\
  S_t &= \max\bigl(0,\; S_{t-1} + v_t - \mu_0 - \kappa\bigr),
  \label{eq:cusum}
\end{align}
where $\kappa > 0$ is a \emph{drift sensitivity} parameter that controls the
trade-off between detection delay and false-alarm rate.  An alarm is raised
when $S_t > h$ for a threshold $h$.
\end{definition}

\begin{theorem}[Spring Detection Guarantees]
\label{thm:spring}
Let $v_t$ be i.i.d.\ $\sigma_0^2$-sub-Gaussian with mean $\mu_0$ under $H_0$
and mean $\mu_1$ under $H_1$.  Set $\kappa = (\mu_1 - \mu_0)/2$.  Then for
any threshold $h > 0$:
\begin{enumerate}[label=(\roman*),nosep]
  \item \textbf{False-alarm rate:}
    $\Pbb(\text{alarm before change} \mid H_0) \le
     \exp\bigl(-\frac{2h\kappa}{\sigma_0^2}\bigr)$;
  \item \textbf{Expected detection delay} (Wald's approximation):
    $\E[T_{\text{detect}} \mid H_1] \le
     \frac{h}{\mu_1 - \mu_0 - \kappa} + \mathcal{O}(1)$.
\end{enumerate}
\end{theorem}

\begin{proof}
Under $H_0$, $v_t - \mu_0$ is zero-mean sub-Gaussian.  The CUSUM process
$S_t$ is a reflected random walk with negative drift $-\kappa$.  By standard
CUSUM analysis~\cite{basseville1993detection,page1954continuous}:
\begin{equation}
  \Pbb(\max_{t \ge 0} S_t > h \mid H_0) \le
  \exp(-2h\kappa / \sigma_0^2),
\end{equation}
which follows from the Chernoff bound applied to the likelihood-ratio
martingale.

For detection delay, under $H_1$ the increment has positive mean
$\mu_1 - \mu_0 - \kappa = \kappa$.  Wald's equation~\cite{wald1945sequential}
gives $\E[T_{\text{detect}}] \approx h/\kappa$, with the $\mathcal{O}(1)$
term accounting for the overshoot.
\end{proof}

\subsection{OOD Detection Lower Bound}
\label{subsec:spring_ood}

The CUSUM statistic above monitors temporal drift across a rollout.  Spring
also induces a one-sample OOD test: a sample is suspicious when its gatekeeper
score is atypical relative to scores stored in the in-distribution memory.  To
avoid collision with the CUSUM statistic $S_t$, write $G_t(x)$ for the
per-sample Spring gatekeeper score; this is the score denoted $S_t(x)$ in the
Spring gating notes.

Let $\mathfrak{M}_t$ be the Spring memory bank of accepted in-distribution
samples.  For $M$ expert modules, let $Z_{m,t}(x) \in [0,1]$ be expert $m$'s
normalized validity score for sample $x$ at time $t$, and define the aggregate
gatekeeper score
\begin{equation}
  G_t(x) = \frac{1}{M}\sum_{m=1}^{M} Z_{m,t}(x).
  \label{eq:spring_gate_score}
\end{equation}
The memory bank estimates the in-distribution score center
$\mu_{\mathrm{in},t} = \E[G_t(X)\mid X \sim P_{\mathrm{in},t}]$ by
$\widehat{\mu}_{\mathrm{in},t}$.

\begin{definition}[Spring OOD Hypothesis Test]
\label{def:spring_ood_test}
OOD detection is the binary hypothesis test
\begin{equation}
  H_0: X \sim P_{\mathrm{in},t},
  \qquad
  H_1: X \sim P_{\mathrm{out},t},
\end{equation}
with OOD shift
\begin{equation}
  \delta_{\mathrm{OOD}}
  =
  \left|
  \E[G_t(X)\mid H_1]
  -
  \mu_{\mathrm{in},t}
  \right|.
\end{equation}
For a tolerance $\epsilon > 0$, the ideal calibrated Spring detector is
\begin{equation}
  \varphi_{\epsilon,t}(x)
  =
  \ind{\left|G_t(x)-\mu_{\mathrm{in},t}\right|>\epsilon},
  \label{eq:spring_ood_test}
\end{equation}
where $\varphi_{\epsilon,t}(x)=1$ means ``declare OOD.''  In deployment,
$\mu_{\mathrm{in},t}$ is replaced by $\widehat{\mu}_{\mathrm{in},t}$ from
$\mathfrak{M}_t$; the calibration error is handled in
Corollary~\ref{cor:spring_memory_scaling}.
\end{definition}

\begin{theorem}[SCX OOD Detection Lower Bound]
\label{thm:spring_ood_lower}
Condition on the Spring memory state $\mathfrak{M}_t$.  Assume that
$Z_{1,t}(X),\ldots,Z_{M,t}(X)$ are independent and take values in $[0,1]$.
Under $H_0$, suppose $\E[G_t(X)\mid H_0]=\mu_{\mathrm{in},t}$; under $H_1$,
suppose
$|\E[G_t(X)\mid H_1]-\mu_{\mathrm{in},t}|=\delta_{\mathrm{OOD}}$.
Then, writing
$p_i(\epsilon)=\Pbb(\varphi_{\epsilon,t}(X)=1\mid H_i)$, the test
in~\eqref{eq:spring_ood_test} satisfies
\begin{align}
  p_0(\epsilon)
    &\le 2\exp(-2M\epsilon^2), \label{eq:ood_false_alarm}\\
  p_1(\epsilon)
    &\ge
    1-\exp\!\left(-2M(\delta_{\mathrm{OOD}}-\epsilon)^2\right)
    \label{eq:ood_power_exact}
\end{align}
for every $0<\epsilon<\delta_{\mathrm{OOD}}$.  In particular, if the OOD
shift has margin $\delta_{\mathrm{OOD}}\ge 2\epsilon$, then
\begin{equation}
  \Pbb(\text{detect OOD}\mid \delta_{\mathrm{OOD}}\ge 2\epsilon)
  \ge 1-\exp(-2M\epsilon^2).
  \label{eq:ood_power_margin}
\end{equation}
\end{theorem}

\begin{proof}
\rigorFull
The false-alarm bound is Hoeffding's inequality applied to the average
$G_t(X)$ of $M$ independent $[0,1]$ random variables:
\begin{equation}
  \Pbb\!\left(
  |G_t(X)-\mu_{\mathrm{in},t}|>\epsilon \mid H_0
  \right)
  \le 2\exp(-2M\epsilon^2).
\end{equation}

For the detection bound, assume first that
$\E[G_t(X)\mid H_1]=\mu_{\mathrm{in},t}+\delta_{\mathrm{OOD}}$; the negative
shift case is identical after reversing the sign.  A missed detection under
$H_1$ implies
\begin{equation}
  G_t(X)-\mu_{\mathrm{in},t}\le \epsilon,
\end{equation}
and therefore
\begin{equation}
  G_t(X)-\E[G_t(X)\mid H_1]
  \le
  -(\delta_{\mathrm{OOD}}-\epsilon).
\end{equation}
Hoeffding's one-sided inequality gives
\begin{equation}
  \Pbb(\varphi_{\epsilon,t}(X)=0\mid H_1)
  \le
  \exp\!\left(-2M(\delta_{\mathrm{OOD}}-\epsilon)^2\right).
\end{equation}
Taking complements proves~\eqref{eq:ood_power_exact}.  If
$\delta_{\mathrm{OOD}}\ge 2\epsilon$, then
$(\delta_{\mathrm{OOD}}-\epsilon)^2\ge \epsilon^2$, yielding
the margin-normalized bound~\eqref{eq:ood_power_margin}.
\end{proof}

\begin{theorem}[Minimum Detectable OOD Shift]
\label{thm:spring_ood_min_shift}
Assume the expert scores are independent with common variance
$\Var(Z_{m,t}(X)\mid H_i)=\sigma^2$ under the local alternatives
$i\in\{0,1\}$.  Then
\begin{equation}
  \Var(G_t(X)\mid H_i)=\frac{\sigma^2}{M}.
\end{equation}
Consequently, the unit-standard-error detectable shift of the Spring
gatekeeper is
\begin{equation}
  \delta_{\min}=\frac{\sigma}{\sqrt{M}}.
  \label{eq:ood_delta_min}
\end{equation}
More generally, in the Gaussian local model
$G_t(X)\mid H_0\sim \cN(\mu_{\mathrm{in},t},\sigma^2/M)$ and
$G_t(X)\mid H_1\sim \cN(\mu_{\mathrm{in},t}+\delta,\sigma^2/M)$, any
level-$\alpha$ one-sided threshold test needs
\begin{equation}
  \delta
  \ge
  \bigl(z_{1-\alpha}+z_{1-\beta}\bigr)\frac{\sigma}{\sqrt{M}}
  \label{eq:ood_delta_alpha_beta}
\end{equation}
to achieve power at least $1-\beta$, where $z_p$ is the $p$-quantile of the
standard normal distribution.
\end{theorem}

\begin{proof}
The variance identity follows from independence:
\begin{align}
  \Var(G_t(X)\mid H_i)
  &=
  \Var\!\left(
  \frac{1}{M}\sum_{m=1}^M Z_{m,t}(X)
  \middle| H_i
  \right) \notag\\
  &=
  \frac{1}{M^2}\sum_{m=1}^M \sigma^2
  =
  \frac{\sigma^2}{M}.
\end{align}
Thus the natural noise scale of the aggregate score is
$\sigma/\sqrt{M}$, proving~\eqref{eq:ood_delta_min}.

For the Gaussian local model, the Neyman--Pearson most powerful level-$\alpha$
one-sided test rejects $H_0$ when
$G_t(X)>\mu_{\mathrm{in},t}+z_{1-\alpha}\sigma/\sqrt{M}$.  Its power under
$H_1$ is
\begin{equation}
  1-\Phi_{\mathrm{std}}\!\left(
  z_{1-\alpha}-\frac{\delta\sqrt{M}}{\sigma}
  \right),
\end{equation}
where $\Phi_{\mathrm{std}}$ is the standard normal CDF.  Requiring this
quantity to be at least $1-\beta$ is equivalent to
$z_{1-\alpha}-\delta\sqrt{M}/\sigma \le z_{\beta}=-z_{1-\beta}$, which is
exactly~\eqref{eq:ood_delta_alpha_beta}.
\end{proof}

\begin{corollary}[Single-Expert Limitation and Memory Scaling]
\label{cor:spring_memory_scaling}
At unit-standard-error resolution, a single Spring expert has
$\delta_{\min}=\sigma$.  Detecting an OOD shift of size $\epsilon$ therefore
requires
\begin{equation}
  M > \frac{\sigma^2}{\epsilon^2}
  \label{eq:ood_expert_requirement}
\end{equation}
experts, up to the confidence constants in
\eqref{eq:ood_delta_alpha_beta}.  If the memory bank contains
$n_t=|\mathfrak{M}_t|$ independent in-distribution calibration samples and
$\widehat{\mu}_{\mathrm{in},t}=n_t^{-1}\sum_{j=1}^{n_t}G_t(X_j)$, then
\begin{equation}
  \Var\!\left(G_t(X)-\widehat{\mu}_{\mathrm{in},t}\mid H_0\right)
  =
  \frac{\sigma^2}{M}\left(1+\frac{1}{n_t}\right),
\end{equation}
and the memory-calibrated detectable shift is
\begin{equation}
  \delta_{\min}(t)
  =
  \frac{\sigma}{\sqrt{M}}
  \sqrt{1+\frac{1}{n_t}}.
  \label{eq:ood_memory_delta_min}
\end{equation}
Thus Spring's OOD sensitivity improves as $\mathfrak{M}_t$ grows, converging
to the expert-ensemble floor $\sigma/\sqrt{M}$.  If the memory mechanism also
increases the effective number of independent reference experts
$M_{\mathrm{eff}}(t)$, replace $M$ by $M_{\mathrm{eff}}(t)$ in
\eqref{eq:ood_memory_delta_min}, giving the monotone improvement
$\delta_{\min}(t)\asymp \sigma/\sqrt{M_{\mathrm{eff}}(t)}$.
\end{corollary}

\begin{proof}
The single-expert statement is Theorem~\ref{thm:spring_ood_min_shift} with
$M=1$.  Solving $\epsilon>\sigma/\sqrt{M}$ for $M$ gives
\eqref{eq:ood_expert_requirement}.  For the memory term, the new-sample score
$G_t(X)$ and the empirical memory mean $\widehat{\mu}_{\mathrm{in},t}$ are
independent under calibration, with variances $\sigma^2/M$ and
$\sigma^2/(Mn_t)$, respectively.  The variance of their difference is the sum
of these two variances, yielding~\eqref{eq:ood_memory_delta_min}.
\end{proof}

\subsection{Multi-Resolution Monitoring}

Spring employs a \emph{multi-resolution} monitoring scheme with three
timescales:

\begin{enumerate}[leftmargin=*,nosep]
  \item \textbf{Fast (per-step):} Conservation-law check (energy, momentum)
    with per-step CUSUM $S_t^{\text{fast}}$;
  \item \textbf{Medium (per-episode):} Trajectory-level drift via
    Wasserstein distance between predicted and expected state distributions,
    $S_e^{\text{med}}$;
  \item \textbf{Slow (per-deployment):} Distributional shift detection via
    maximum mean discrepancy (MMD) on the full deployment history,
    $S_d^{\text{slow}}$.
\end{enumerate}

Alarms are raised via a hierarchical voting scheme: a fast alarm triggers a
medium-resolution check; a medium alarm triggers a slow-resolution audit.
This cascaded design balances sensitivity with false-alarm control.

\begin{algorithm}[t]
\caption{Spring Online Drift Monitor}
\label{alg:spring}
\begin{algorithmic}[1]
\Require World model $M$, aggregate certificate $\psi_{\text{agg}}$,
         thresholds $(h_f, h_m, h_s)$, reference distribution $P_0$
\Ensure Alarms at three resolutions
\State Initialize $S_f \gets 0$, $S_m \gets 0$, $S_s \gets 0$, $t \gets 0$
\For{each deployment step}
  \State $t \gets t + 1$
  \State $\hat{\mathbf{s}} \gets M(\mathbf{s}_{t-1}, \mathbf{x}_{t-1})$
  \State $v_t \gets -\log \psi_{\text{agg}}(\hat{\mathbf{s}})$
  \State $S_f \gets \max(0, S_f + v_t - \mu_0 - \kappa_f)$
  \If{$S_f > h_f$}
    \State \textbf{emit} FAST\_ALARM
    \State Compute $\text{MMD}(\hat{P}_t^{\text{episode}}, P_0)$
    \State $S_m \gets S_m + \text{MMD}$
    \If{$S_m > h_m$} \textbf{emit} MEDIUM\_ALARM \EndIf
  \EndIf
  \If{episode ends}
    \State Update slow detector: $S_s \gets \text{MMD}(\hat{P}_{\text{deployment}}, P_0)$
    \If{$S_s > h_s$} \textbf{emit} SLOW\_ALARM \EndIf
  \EndIf
\EndFor
\end{algorithmic}
\end{algorithm}

\subsection{Calibration Phase}

Before deployment, Spring undergoes a \emph{calibration phase} on a held-out
validation set of physically valid trajectories.  During calibration:
\begin{enumerate}[leftmargin=*,nosep]
  \item Estimate $\mu_0$, $\sigma_0^2$ from empirical violation scores;
  \item Set thresholds $(h_f, h_m, h_s)$ via permutation testing to achieve a
    target false-alarm rate $\alpha$ (e.g., $\alpha = 0.01$ per deployment hour);
  \item Validate on adversarial perturbations to ensure sensitivity to known
    failure modes.
\end{enumerate}

%==============================================================================
\section{Cercis Score for Learned Simulators}
\label{sec:cercis}
%==============================================================================

The \textbf{Cercis Score} is a composite metric designed to rate the
trustworthiness of a learned world model for a given deployment scenario.  It
combines two orthogonal dimensions:

\begin{definition}[Cercis Score]
For world model $M$ and deployment scenario $\sigma$, the Cercis Score is:
\begin{equation}\label{eq:cercis}
  C(M, \sigma) = Q(M, \sigma) \cdot N(\sigma),
\end{equation}
where:
\begin{itemize}[nosep]
  \item $Q(M, \sigma) \in [0,1]$: \textbf{Physical Fidelity}---quantifies how
    well $M$ respects physical invariants under scenario $\sigma$;
  \item $N(\sigma) \in [0,1]$: \textbf{Scenario Novelty}---quantifies how
    out-of-distribution scenario $\sigma$ is relative to the model's training
    manifold.
\end{itemize}
\end{definition}

\subsection{Physical Fidelity $Q$}

$Q$ aggregates multiple physics-invariant checks into a single scalar:

\begin{equation}
  Q(M, \sigma) =
  \frac{1}{|\mathcal{I}|} \sum_{\mathcal{I}_j \in \mathcal{I}}
  w_j \cdot q_j(M, \sigma),
\end{equation}
where $\mathcal{I} = \{\mathcal{I}_E, \mathcal{I}_p, \mathcal{I}_{\cap},
\mathcal{I}_f, \dots\}$ is the set of audited invariants, $w_j$ are
domain-specific weights, and $q_j$ is the per-invariant fidelity score.

\paragraph{Conservation Law Fidelity.}
For energy conservation ($\mathcal{I}_E$), we measure the normalized energy
drift over a rollout of length $T$:
\begin{equation}
  q_E(M, \sigma) =
  \exp\!\left(
    -\lambda_E \cdot \frac{1}{T}\sum_{t=1}^{T}
    \frac{|E_t - E_0 - W_{0:t}|}{|E_0| + \varepsilon}
  \right),
\end{equation}
where $E_t$ is the system energy at step $t$, $W_{0:t}$ is the cumulative
external work, and $\lambda_E$ controls sensitivity.  Analogous definitions
hold for linear momentum ($q_p$) and angular momentum ($q_L$).

\paragraph{Collision Detection Fidelity.}
For non-penetration ($\mathcal{I}_{\cap}$):
\begin{equation}
  q_{\cap}(M, \sigma) =
  \frac{1}{T\cdot |\text{pairs}|}
  \sum_{t=1}^{T} \sum_{(i,j)}
  \ind{\phi_t(\mathbf{q}_i, \mathbf{q}_j) \ge -\varepsilon_{\cap}},
\end{equation}
where $\phi_t$ is the signed distance at step $t$ and $\varepsilon_{\cap}$ is
a tolerance (typically $10^{-6}$ m).

\paragraph{Friction Cone Fidelity.}
For the friction cone invariant ($\mathcal{I}_f$):
\begin{equation}
  q_f(M, \sigma) =
  \frac{1}{T\cdot |\text{contacts}|}
  \sum_{t=1}^{T} \sum_{c}
  \ind{\norm{\mathbf{f}_t^{(c)}} \le \mu_c \norm{\mathbf{f}_n^{(c)}} + \varepsilon_f}.
\end{equation}

\subsection{Scenario Novelty $N$}

$N(\sigma)$ measures distributional deviation from the training manifold:

\begin{equation}
  N(\sigma) = \min\!\left(1,\;
    \frac{\text{MMD}(\cD_\sigma, \cD_{\text{train}})}
         {\text{MMD}_{\text{ref}}}
  \right),
\end{equation}
where $\text{MMD}$ is the maximum mean discrepancy with an RBF kernel,
$\cD_\sigma$ is the empirical distribution of states visited under scenario
$\sigma$, $\cD_{\text{train}}$ is the training distribution, and
$\text{MMD}_{\text{ref}}$ is a reference value (e.g., the 95th percentile of
in-distribution MMD values).

\begin{remark}[Interpretation]
\begin{itemize}[nosep]
  \item $C \approx 1$: High physical fidelity in familiar territory---the
    model is trustworthy for this scenario;
  \item $C \ll 1$ with $Q \approx 1, N \ll 1$: Good physics but routine
    scenario (low information value for auditing);
  \item $C \ll 1$ with $Q \ll 1, N \approx 1$: Novel scenario exposing
    physics violations---\textbf{high-risk, requires human review};
  \item $C \ll 1$ with $Q \ll 1, N \ll 1$: Fundamental physics failures even
    in familiar settings---\textbf{model is broken}.
\end{itemize}
\end{remark}

\subsection{Connection to Theorem~3}

The Cercis Score directly addresses the unidentifiability result of
Theorem~\ref{thm:unident} by making the untestable assumptions \emph{explicit
and quantifiable}: $Q$ measures what the model claims (physical fidelity),
$N$ measures how far the deployment deviates from the assumptions under which
$Q$ was calibrated.  A Cercis Score $C > \tau$ (for a domain-specific
threshold $\tau$) translates to the operational guarantee:

\begin{quote}
  \emph{Under the declared assumptions (embedded in the training distribution
  $\cD_{\text{train}}$ and invariant set $\mathcal{I}$), the model's physical
  fidelity is at least $Q$ with novelty bounded by $N$.}
\end{quote}

This conditional guarantee is the strongest possible statement given
Theorem~\ref{thm:unident}.

%==============================================================================
\section{Multi-Module Yajie Consensus}
\label{sec:yajie}
%==============================================================================

When a world model comprises $K$ heterogeneous modules, audit signals from
individual modules may conflict.  The \textbf{Yajie Consensus} protocol
(Yajie Consensus) resolves these conflicts through a Byzantine-resilient
aggregation mechanism.

\subsection{Problem Formulation}

Each module $k$ produces an audit signal $a_k \in [0,1]$ indicating its
confidence that the composite output is physically valid.  However:
\begin{itemize}[nosep]
  \item Some modules may be \emph{overconfident} (neural modules with poor
    calibration);
  \item Some may be \emph{underconfident} (conservative physics engines);
  \item In adversarial settings, up to $f$ modules may be \emph{Byzantine}
    (arbitrarily malicious audit signals).
\end{itemize}

The goal is to produce a consensus audit signal $\bar{a}$ that is robust to
up to $f$ Byzantine modules, assuming the remaining $K - f$ modules are
\emph{honest} (their audit signals are well-calibrated).

\subsection{Trimmed-Mean Consensus}

\begin{definition}[Yajie Trimmed-Mean Operator]
Given module audit signals $\{a_k\}_{k=1}^K$, the Yajie consensus signal is:
\begin{equation}\label{eq:yajie}
  \bar{a} = \text{TrimMean}_f(\{a_k\}) =
  \frac{1}{K - 2f}
  \sum_{k=f+1}^{K-f} a_{(k)},
\end{equation}
where $a_{(1)} \le a_{(2)} \le \cdots \le a_{(K)}$ are the sorted audit
signals.  The $f$ smallest and $f$ largest values are discarded.
\end{definition}

\begin{theorem}[Yajie Robustness]
\label{thm:yajie}
Assume $K > 3f$ and that honest modules produce audit signals
$a_k^{\text{honest}}$ satisfying $|a_k^{\text{honest}} - a^*| \le \delta$ for
some ground-truth physical validity $a^* \in [0,1]$.  Then the Yajie
consensus satisfies:
\begin{equation}
  |\bar{a} - a^*| \le \delta + \frac{2f}{K-2f}.
\end{equation}
\end{theorem}

\begin{proof}
After trimming the $f$ smallest and $f$ largest values, at most $f$ Byzantine
signals remain among the $K-2f$ retained values.  In the worst case, all $f$
retained Byzantine signals are at the extremes ($0$ or $1$), contributing at
most $\frac{2f}{K-2f}$ to the deviation.  The remaining $K-3f$ honest signals
each deviate from $a^*$ by at most $\delta$.  The trimmed mean bounds follow
by direct computation:

\begin{align}
  \bar{a} &\le \frac{(K-3f)(a^* + \delta) + f \cdot 1}{K-2f}
            = a^* + \delta + \frac{f(1 - a^* - \delta)}{K-2f} \\
          &\le a^* + \delta + \frac{f}{K-2f},
\end{align}
and symmetrically for the lower bound.
\end{proof}

\subsection{Weighted Yajie with Cercis Priors}

In practice, we augment the trimmed mean with module-specific weights derived
from Cercis calibration:

\begin{equation}
  \bar{a}_{\text{weighted}} =
  \frac{\sum_{k=f+1}^{K-f} w_{(k)} \, a_{(k)}}
       {\sum_{k=f+1}^{K-f} w_{(k)}},
\end{equation}
where $w_k = Q_k \cdot \exp(-\gamma \cdot \text{calibration\_error}_k)$,
with $Q_k$ being the module's per-invariant Cercis $Q$-score and
$\text{calibration\_error}_k$ measuring the discrepancy between module $k$'s
certified confidence and its empirical error rate on a validation set.

\begin{algorithm}[t]
\caption{Weighted Yajie Consensus}
\label{alg:yajie}
\begin{algorithmic}[1]
\Require Audit signals $\{a_k\}_{k=1}^K$, weights $\{w_k\}_{k=1}^K$,
         Byzantine tolerance $f$
\Ensure Consensus signal $\bar{a}$, flag for potential Byzantine modules
\State Sort pairs $\{(a_k, w_k)\}$ by $a_k$ ascending
\State Trim $f$ smallest and $f$ largest
\State $\bar{a} \gets \sum_{k=f+1}^{K-f} w_{(k)} a_{(k)} \,/\, \sum_{k=f+1}^{K-f} w_{(k)}$
\For{$k \in \{1,\dots,f\} \cup \{K-f+1,\dots,K\}$}
  \If{$|a_k - \bar{a}| > 3\delta$}
    \State Flag module $k$ as potentially Byzantine
  \EndIf
\EndFor
\State \Return $\bar{a}$, flagged modules
\end{algorithmic}
\end{algorithm}

\subsection{Application to Isaac Sim Pipelines}

In an Isaac Sim pipeline with five modules (scene graph, PhysX, neural
residual, RTX renderer, Cosmos prior), the Yajie protocol operates as
follows:
\begin{enumerate}[leftmargin=*,nosep]
  \item Each module computes its per-invariant audit signals on its own output
    (e.g., PhysX reports collision-penetration depth, Cosmos reports
    diffusion-loss statistics);
  \item Audit signals are collected and the weighted trimmed mean is computed
    with $f=1$ (tolerating one potentially Byzantine module);
  \item If consensus $\bar{a} < \tau_{\text{safe}}$, the pipeline is paused
    for human review;
  \item Modules persistently flagged as Byzantine trigger a
    recalibration/replacement workflow.
\end{enumerate}

\subsection{Convergence Analysis and Computational Overhead}

We briefly analyze the convergence properties and computational cost of the
Yajie protocol.

\begin{proposition}[Consensus Convergence Rate]
\label{prop:yajie_convergence}
Under the honest-module assumption ($|a_k^{\text{honest}} - a^*| \le \delta$),
the Yajie weighted consensus $\bar{a}_t$ at round $t$ of an iterative audit
loop converges to within $\delta$ of $a^*$ at a geometric rate:
\begin{equation}
  \E[|\bar{a}_t - a^*|] \le \delta + \frac{2f}{K-2f} +
  \left(1 - \frac{K-3f}{K-2f}\right)^t \cdot \Delta_0,
\end{equation}
where $\Delta_0 = |\bar{a}_0 - a^*|$ is the initial consensus error.
\end{proposition}

\begin{proof}
Each round retains $K-2f$ modules, of which at least $K-3f$ are honest.
The contribution of Byzantine signals is bounded by $2f/(K-2f)$, while the
honest signals' error contracts by a factor of $1 - (K-3f)/(K-2f) = f/(K-2f)$
per round due to the weighted averaging with recalibrated weights.
\end{proof}

\paragraph{Computational Overhead.}
The Yajie protocol adds negligible overhead to the simulation pipeline:
\begin{itemize}[nosep]
  \item Sorting $K$ audit signals: $\mathcal{O}(K \log K)$ per audit cycle;
  \item Weighted mean computation: $\mathcal{O}(K)$ per cycle;
  \item Byzantine flagging: $\mathcal{O}(K)$ per cycle.
\end{itemize}
For typical $K \in [3, 10]$, this is sub-microsecond on modern hardware---well
below the per-step simulation cost (PhysX integration, neural-network forward
passes).  The dominant cost is the per-module certificate computation
$\tilde{\psi}_k$, which can be amortized across multiple audit cycles or
computed asynchronously.

%==============================================================================
\section{Experimental Protocol}
\label{sec:experiments}
%==============================================================================

We specify a comprehensive experimental protocol for validating the SCX
framework on NVIDIA Isaac Sim~\cite{isaacsim2023} benchmarks with real-robot
transfer.

\subsection{Benchmark Suite}

\begin{table}[t]
\centering
\caption{SCX World-Model Audit Benchmark Suite}
\label{tab:benchmarks}
\begin{tabular}{@{}llc@{}}
\toprule
\textbf{Task} & \textbf{Domain} & \textbf{Modules} \\
\midrule
Franka Cube Stacking   & Manipulation  & 5 \\
ANYmal Locomotion      & Locomotion    & 5 \\
Shadow Hand In-Hand    & Dexterity     & 5 \\
Drone Obstacle Course  & Aerial        & 4 \\
Soft-body Grasping     & Deformable    & 5 \\
\bottomrule
\end{tabular}
\end{table}

\subsection{Evaluation Metrics}

\begin{enumerate}[leftmargin=*,nosep]
  \item \textbf{Cascade Miss Rate} ($\hat{\beta}_{\text{cascade}}$):
    Empirical probability that a physics violation introduced at any module
    survives undetected through all subsequent detectors.  Compared against
    the theoretical bound of Theorem~\ref{thm:cascade}.
  \item \textbf{Spring Detection Delay}: Time (in simulation steps) from
    violation onset to Spring alarm, at controlled false-alarm rates
    $\alpha \in \{0.001, 0.01, 0.05\}$.
  \item \textbf{Cercis Score Correlation}: Correlation between $C(M,\sigma)$
    and real-robot task success rate across 50 deployment scenarios with
    varying novelty.
  \item \textbf{Assumption Violation Rate}: For Theorem~\ref{thm:unident},
    measure how often sim-to-real transfer pipelines that do \emph{not}
    declare assumptions silently fail vs.\ those using SCX assumption
    declarations.
\end{enumerate}

\subsection{Protocol Details}

\paragraph{Phase 1: Offline Audit (Isaac Sim).}
\begin{enumerate}[leftmargin=*,nosep]
  \item \textbf{Calibration}: Run 10,000 rollouts of each benchmark on Isaac
    Sim.  Compute per-module certificate calibrations
    $\{\tilde{\psi}_k\}$, estimate Lipschitz constants $\{L_k\}$ via
    finite-difference perturbation, and calibrate Spring thresholds.
  \item \textbf{Violation Injection}: Introduce controlled physics violations
    at each module (e.g., zero a friction coefficient, disable collision
    between a specific object pair, perturb a neural module's latent code).
    Measure cascade miss rates and compare against the union bound and cascade
    bound.
  \item \textbf{Adversarial Audit}: Apply gradient-based adversarial attacks
    to each neural module (PGD on the dynamics residual, adversarial
    perturbations to the Cosmos diffusion prior).  Verify Spring detection
    sensitivity.
\end{enumerate}

\paragraph{Phase 2: Real-Robot Transfer.}
\begin{enumerate}[leftmargin=*,nosep]
  \item \textbf{Policy Training}: Train RL policies (PPO~\cite{schulman2017ppo})
    for each task entirely in Isaac Sim with domain randomization.
  \item \textbf{Sim-to-Real Deployment}: Deploy policies on physical hardware
    (Franka Emika Panda for manipulation, ANYmal C for locomotion).
    Simultaneously run Spring monitoring on a shadow simulation.
  \item \textbf{Cercis Scoring}: For each deployment scenario, compute
    $C(M,\sigma)$.  Record per-episode success/failure.
  \item \textbf{Assumption Tracking}: For each deployment, document which
    structural assumptions were used (e.g., ``friction coefficient $\mu \in
    [0.3, 0.8]$'', ``object mass $\pm 10\%$'').  Correlate assumption
    violations with deployment failures.
\end{enumerate}

\paragraph{Phase 3: Ablation Studies.}
\begin{enumerate}[leftmargin=*,nosep]
  \item \textbf{Number of Modules}: Vary $K \in \{2,3,4,5\}$ by
    enabling/disabling modules.  Measure cascade bound tightness.
  \item \textbf{Byzantine Tolerance}: Inject $f \in \{0,1,2\}$ Byzantine
    modules with malicious audit signals.  Compare Yajie consensus against
    simple mean and median aggregation.
  \item \textbf{Novelty Gradient}: Vary scenario novelty $N$ from 0 to 1 via
    controlled OOD perturbations (novel object geometries, extreme friction
    values, unseen lighting conditions).  Track $Q$ degradation.
\end{enumerate}

\subsection{Expected Results}

Based on preliminary analysis and Theorem~\ref{thm:cascade}, we expect:

\begin{enumerate}[leftmargin=*,nosep]
  \item The cascade bound (Eq.~\ref{eq:cascade_bound}) to be $3$--$10\times$
    tighter than the union bound for $K \ge 4$ modules, with $\beta_k \in
    [0.05, 0.20]$;
  \item Spring to detect 95\% of injected violations within 50 simulation
    steps at $\alpha = 0.01$ false-alarm rate;
  \item Cercis Score $C$ to correlate with real-robot success rate at
    $\rho > 0.75$ (Spearman);
  \item Yajie consensus to reject Byzantine audit signals with less than 5\%
    degradation in consensus accuracy when $f \le \lfloor (K-1)/3 \rfloor$;
  \item Sim-to-real pipelines without declared assumptions to exhibit
    unpredictable failure modes in $\ge 30\%$ of novel scenarios, consistent
    with Theorem~\ref{thm:unident}.
\end{enumerate}

\subsection{Statistical Power and Sample-Size Analysis}

To ensure the experimental protocol has sufficient statistical power, we
derive minimum sample sizes for the key hypothesis tests.

\paragraph{Cascade Bound Tightness.}
Let $\hat{r} = \hat{\beta}_{\text{union}} / \hat{\beta}_{\text{cascade}}$ be
the empirical tightness ratio.  To detect a ratio $r > 3$ with power
$1-\gamma = 0.9$ at significance level $\alpha = 0.05$, the per-module sample
size $n$ must satisfy:
\begin{equation}
  n \ge \frac{2(z_{\alpha/2} + z_\gamma)^2 \sigma_r^2}{(\log r)^2},
\end{equation}
where $\sigma_r^2$ is the variance of $\log \hat{r}$ estimated from pilot
experiments, and $z_p$ denotes the $p$-quantile of the standard normal
distribution.  Pilot experiments with $K=4$, $\beta_k = 0.1$ yield
$\sigma_r^2 \approx 0.15$, requiring $n \ge 420$ violation-injection trials
per module.

\paragraph{Cercis-Task-Success Correlation.}
To detect a Spearman correlation $\rho \ge 0.75$ with power 0.9, the required
number of deployment scenarios is:
\begin{equation}
  N_{\text{scenarios}} \ge
  \left(\frac{z_{\alpha/2} + z_\gamma}{\operatorname{arctanh}(\rho)}\right)^2 + 3.
\end{equation}
For $\rho = 0.75$, $\alpha = 0.05$, $\gamma = 0.1$, this yields
$N_{\text{scenarios}} \ge 13$.  We conservatively use 50 scenarios to
enable subgroup analysis (by task type, novelty level, and module
configuration).

\paragraph{Spring Detection Sensitivity.}
Under the CUSUM detection model (Theorem~\ref{thm:spring}), the expected
detection delay $\E[T_{\text{detect}}]$ scales inversely with the drift
magnitude $\mu_1 - \mu_0$.  To achieve $\E[T_{\text{detect}}] \le 50$ steps
with 95\% confidence:
\begin{equation}
  \mu_1 - \mu_0 \ge \frac{h}{50} + \kappa,
\end{equation}
where $h$ is calibrated from the false-alarm constraint.  For $\alpha = 0.01$
and $\sigma_0^2 = 0.25$ (estimated from calibration), this requires a drift
magnitude of at least $\mu_1 - \mu_0 \ge 0.15$---achievable for the
violation types in our benchmark (e.g., disabling collision detection raises
$v_t$ by $0.3$--$0.7$).  For one-sample OOD detection, the complementary
Spring gatekeeper bound in Theorem~\ref{thm:spring_ood_lower} predicts
exponential power growth in the number of independent experts, while
Corollary~\ref{cor:spring_memory_scaling} gives the detectable shift scale
$\delta_{\min}(t)$ after memory calibration.

%==============================================================================
\section{Discussion}
\label{sec:discussion}
%==============================================================================

\subsection{Summary}

We have presented the SCX audit framework for learned world models
(world models), establishing formal guarantees (Theorems~\ref{thm:cascade}
and~\ref{thm:unident}), runtime monitoring (Spring), a composite trust metric
(Cercis Score), and a Byzantine-resilient consensus protocol (Yajie Consensus).  The framework addresses a critical gap in the current
world-model ecosystem: while systems like NVIDIA Cosmos and Isaac Sim achieve
impressive perceptual and predictive quality, their \emph{physical
consistency} remains uncertified, creating safety risks for downstream
applications in robotics, autonomous driving, and scientific simulation.

\subsection{Limitations}

\begin{enumerate}[leftmargin=*,nosep]
  \item \textbf{Certificate Conservatism}: Conservative certificates
    $\tilde{\psi}_k$ may significantly underestimate true module fidelity,
    weakening the cascade bound.  Adaptive certificate calibration
    (e.g., via conformal prediction with online updates) is an important
    direction.
  \item \textbf{Lipschitz Estimation}: Module Lipschitz constants $L_k$ are
    challenging to estimate for black-box neural components.  Spectral
    normalization~\cite{miyato2018spectral} and Lipschitz-constrained
    architectures~\cite{trockman2025mim} can provide certified bounds but may
    reduce model capacity.
  \item \textbf{Byzantine Model}: Our Yajie analysis assumes at most $f < K/3$
    Byzantine modules, which may be optimistic for highly heterogeneous
    pipelines.  Broader Byzantine models (e.g., $f < K/2$ with cryptographic
    commitments) remain an open challenge.
  \item \textbf{Unidentifiability Scope}: Theorem~\ref{thm:unident} applies to
    the observational setting; additional structure (e.g., known parametric
    forms for the domain gap, access to counterfactual
    interventions~\cite{pearl2009causality}) may restore identifiability,
    consistent with Corollary~\ref{cor:assumption}.
\end{enumerate}

\subsection{Relationship to Existing Work}

Our work intersects several active research areas:

\paragraph{Simulation Fidelity.}
Traditional verification and validation (V\&V)~\cite{oberkampf2010verification}
focuses on grid-convergence and code verification for classical PDE solvers.
SCX extends V\&V principles to \emph{learned} components with formal
probabilistic guarantees.

\paragraph{World Model Auditing.}
Concurrent work on world-model evaluation~\cite{hu2024eval,michel2025genie}
emphasizes perceptual quality and next-frame prediction accuracy.  SCX
complements these with \emph{physics-specific} invariants.

\paragraph{Out-of-Distribution Detection.}
The novelty component $N$ of the Cercis Score relates to OOD detection
literature~\cite{hendrycks2019baseline,yang2022ood}.  However, $N$ is
specialized to \emph{physical} OOD: states that violate physical invariants
are not merely distributionally novel but physically impossible.

\paragraph{Sim-to-Real Transfer} (\textbf{sim-to-real}~).
Theorem~\ref{thm:unident} provides a theoretical foundation for the empirical
observation that sim-to-real transfer is brittle~\cite{zhao2020sim}.
Our Assumption Declaration protocol offers a path toward
\emph{assumption-aware} transfer.

\paragraph{Formal Verification of Learned Components.}
The formal methods community has developed techniques for verifying neural
network properties including Lipschitz bounds~\cite{fazlyab2019efficient},
adversarial robustness~\cite{zhang2018crown,wang2021beta}, and
safety~\cite{katz2017reluplex}.  SCX adapts these tools to the
\emph{compositional} setting: rather than verifying a single neural network
in isolation, we verify the physical consistency of a cascade of
heterogeneous modules (physics engines, neural networks, renderers).  The
cascade bound (Theorem~\ref{thm:cascade}) is novel in that it accounts for
\emph{inter-module error propagation}---a phenomenon absent from
single-network verification but central to composed world models.

\paragraph{Byzantine-Resilient Aggregation.}
The Yajie protocol draws on classical results in Byzantine fault
tolerance~\cite{lamport1982byzantine,blanchard2017machine} and robust
statistics.  Unlike standard gradient or parameter aggregation, Yajie
aggregates \emph{audit signals}---probabilistic assessments of physical
validity---which have different statistical properties
(bounded support, calibration semantics).  The trimmed-mean approach we
adopt is closely related to the Krum and median-of-means literature but
specialized to the audit-signal domain.

\subsection{Future Directions}

\begin{enumerate}[leftmargin=*,nosep]
  \item \textbf{Neural Certificate Synthesis}: Automatically synthesize
    conservative certificates $\tilde{\psi}_k$ for arbitrary neural modules
    via convex relaxation (e.g., CROWN~\cite{zhang2018crown},
    $\alpha,\beta$-CROWN~\cite{wang2021beta}) or abstraction-refinement loops.
    Recent advances in neural Lyapunov functions~\cite{chang2019neural} and
    barrier certificates~\cite{qin2021safe} for dynamical systems may provide
    template-based synthesis for learned dynamics modules.
  \item \textbf{Differentiable Cercis Score}: Make $C$ differentiable with
    respect to model parameters, enabling \emph{Cercis-regularized training}
    that directly optimizes for physical fidelity.  This could be realized
    through differentiable physics invariants (e.g., automatic differentiation
    of Lagrangian mechanics constraints) integrated into the world-model loss
    function.
  \item \textbf{Federated Spring}: Extend Spring to federated deployments
    where multiple instances of the same world model operate in different
    environments, sharing drift statistics via secure aggregation.  This
    enables fleet-level anomaly detection while preserving operational
    privacy of individual robot deployments.
  \item \textbf{Causal Domain-Gap Identification}: Leverage causal
    discovery~\cite{spirtes2000causation} and interventional data to provide
    stronger identifiability guarantees than Theorem~\ref{thm:unident},
    conditional on the causal graph being known.  In particular, if the causal
    structure of the sim-to-real shift can be partially identified (e.g., which
    environment variables are shifted and by what mechanism), the
    decomposition in Theorem~\ref{thm:unident} becomes identifiable.
  \item \textbf{Hardware-in-the-Loop SCX}: Deploy SCX auditing directly on
    edge compute (Jetson Orin) co-located with robots, enabling real-time
    Cercis scoring and Spring monitoring with sub-millisecond latency.
    Early-stage integration with ROS~2 middleware would allow
    SCX-as-a-service for the robotics ecosystem.
  \item \textbf{Multi-Physics Extensions}: Extend the SCX invariant set beyond
    rigid-body mechanics to include fluid dynamics (Navier--Stokes
    conservation), electromagnetic constraints (Maxwell's equations in
    differentiable form), and thermodynamic laws (entropy production
    monitors).  Each domain introduces new invariant structures and
    module-specific certificates.
  \item \textbf{Uncertainty-Aware Cercis Score}: Incorporate epistemic and
    aleatoric uncertainty estimates into the Cercis decomposition.  When
    $Q$ is low, distinguish between reducible uncertainty (model can be
    improved with more data) and irreducible uncertainty (inherent stochasticity
    in the physical process).
\end{enumerate}

\subsection{Broader Impact and Sociotechnical Considerations}

The SCX framework sits at the intersection of formal methods, machine
learning, and safety engineering.  Its deployment raises several
sociotechnical considerations:

\paragraph{Safety-Critical Deployment.}
World models are increasingly deployed in safety-critical domains: autonomous
driving (Waymo~\cite{waymo2023}, Cruise), surgical robotics (da Vinci
simulation), and nuclear facility inspection.  In these settings, an
undetected physics violation can cause catastrophic harm.  SCX auditing
provides a \emph{necessary but not sufficient} safety layer: a clean Cercis
Score does not guarantee safety, but a poor Cercis Score constitutes a
warning that human operators must not ignore.  We advocate for regulatory
frameworks that mandate assumption declarations and periodic SCX audits for
world-model-based control systems, analogous to the FDA's software-as-medical-device
guidelines.

\paragraph{Over-Reliance and Automation Bias.}
A high Cercis Score might induce unwarranted trust---an operator seeing
$C \approx 1$ may assume the world model is infallible.  The framework must
be deployed with clear communication that $C$ is \emph{conditional} on the
declared assumptions (Corollary~\ref{cor:assumption}) and that physical
reality can always present scenarios outside any finite set of invariants.
We recommend that SCX audit reports always include (i) the full set of
checked invariants, (ii) the declared assumptions, and (iii) the Spring
false-alarm and miss-rate calibration data.

\paragraph{Benchmarking and Reproducibility.}
The Cercis Score provides a standardized metric for comparing world models.
We encourage the community to adopt it alongside perceptual metrics (FVD,
PSNR, LPIPS) so that physical fidelity is tracked with the same rigor as
visual quality.  Open-source implementations of Spring and Yajie will be
released under the SCX project to facilitate reproducible benchmarking
across simulators (Isaac Sim, MuJoCo, Bullet, SAPIEN).

\paragraph{Chinese Simulation Ecosystem}
(Chinese Simulation Ecosystem).
The rapid growth of domestic world-model and simulation platforms in the
Chinese AI ecosystem---including products from SenseTime, Baidu Apollo
simulation, and Horizon Robotics' simulation stack---creates an urgent need
for cross-platform physical consistency standards.  SCX provides a
language-agnostic mathematical framework that can be adopted by any
simulation pipeline regardless of implementation language or geographic
origin.  We explicitly include Chinese-domain terminology
(world models, physical simulation, simulation-to-reality,
Yajie Consensus) to facilitate adoption in the broader international
simulation-verification community.

\subsection{Conclusion}

As learned world models become the backbone of
autonomous systems, their physical consistency guarantees must match their
perceptual quality.  The SCX framework provides a principled foundation:
formal bounds on cross-module violation detection, a proof of fundamental
unidentifiability that demands assumption transparency, and practical tools
(Spring, Cercis, Yajie) for runtime auditing.  We hope this work catalyzes a
shift from \emph{best-effort} simulation to \emph{certified physical
simulation} (\textbf{}), where every sim-to-real transfer
(\textbf{sim-to-real}) is accompanied by a verifiable Cercis Score and an
explicit Assumption Declaration.

%==============================================================================
%  APPENDIX: Extended Proofs
%==============================================================================
\appendix

\section{Extended Proof of Theorem~\ref{thm:cascade}}

\subsection{Formal Derivation of the Lipschitz Amplification Factor}

We provide a rigorous derivation of the amplification factor
$\frac{L_k \cdot \norm{\Delta\mathbf{s}^{(k-1)}}}
       {\tilde{\psi}_k(\mathbf{s}^{(k-1)}, \mathbf{x}_k) + \varepsilon}$
used in the cascade bound.

\begin{lemma}[Certificate Sensitivity Under Input Perturbation]
\label{lem:cert_sensitivity}
Let $M_k$ be $L_k$-Lipschitz with respect to the $\ell_2$ norm:
$\norm{M_k(\mathbf{s}; \mathbf{x}) - M_k(\mathbf{s}'; \mathbf{x})} \le
L_k \norm{\mathbf{s} - \mathbf{s}'}$ for all $\mathbf{s}, \mathbf{s}'$.
Let $\tilde{\psi}_k$ be a conservative certificate: for any input
$\mathbf{s}$, $\tilde{\psi}_k(\mathbf{s}, \mathbf{x}) \le
\Pbb(M_k(\mathbf{s}; \mathbf{x}) \in \mathcal{I}_k \mid \text{module }k\text{'s model})$.

Then for any input $\mathbf{s}$ with nearest valid state $\mathbf{s}_*$:
\begin{equation}
  \Pbb(D_k = 0 \mid \mathbf{s} \text{ invalid}) \le
  \frac{L_k \norm{\mathbf{s} - \mathbf{s}_*}}
       {\tilde{\psi}_k(\mathbf{s}_*, \mathbf{x}) + \varepsilon}.
\end{equation}
\end{lemma}

\begin{proof}
Let $f_k(\mathbf{s}) = \Pbb(M_k(\mathbf{s}) \in \mathcal{I}_k)$ be the true
validity probability.  By the Lipschitz property of $M_k$, for any
$\mathbf{s}, \mathbf{s}'$:
\begin{align}
  |f_k(\mathbf{s}) - f_k(\mathbf{s}')|
  &\le \TV\bigl( M_k(\mathbf{s}), M_k(\mathbf{s}') \bigr) \\
  &\le L_k \norm{\mathbf{s} - \mathbf{s}'},
\end{align}
where the first inequality follows from the data-processing inequality for
total variation and the second from the Lipschitz continuity of $M_k$.

For an invalid input $\mathbf{s}$, let $\mathbf{s}_* = \argmin_{\mathbf{s}'
\in \mathcal{I}_k \text{ (preimage)}} \norm{\mathbf{s} - \mathbf{s}'}$ be the
nearest valid preimage.  Then:
\begin{align}
  \Pbb(M_k(\mathbf{s}) \in \mathcal{I}_k)
  &\le \Pbb(M_k(\mathbf{s}_*) \in \mathcal{I}_k) + L_k \norm{\mathbf{s} - \mathbf{s}_*} \\
  &\le \tilde{\psi}_k(\mathbf{s}_*) + L_k \norm{\Delta\mathbf{s}}.
\end{align}

The detector miss probability is the complement of validity, bounded above by
$1 - \tilde{\psi}_k(\mathbf{s}_*) - L_k \norm{\Delta\mathbf{s}}$.  Using
$\frac{1}{x+\varepsilon} \ge 1-x$ for $x \in [0,1]$ and small $\varepsilon$,
we obtain the stated bound.
\end{proof}

\subsection{High-Probability Bound via Hoeffding}

\begin{lemma}[Finite-Sample Concentration]
\label{lem:hoeffding}
Given $N$ i.i.d.\ calibration samples per module, let $\hat{\beta}_k$ be the
empirical miss rate.  Then with probability at least $1-\delta/K$:
\begin{equation}
  |\hat{\beta}_k - \beta_k| \le \sqrt{\frac{\log(2K/\delta)}{2N}}.
\end{equation}
\end{lemma}

\begin{proof}
Direct application of Hoeffding's inequality to $N$ Bernoulli trials with
success probability $\beta_k$, union bound over $K$ modules.
\end{proof}

Combining Lemmas~\ref{lem:cert_sensitivity} and~\ref{lem:hoeffding} with the
induction argument in the main text yields the $1-\delta$ confidence statement
of Theorem~\ref{thm:cascade}.  $\square$

\section{Extended Proof of Theorem~\ref{thm:unident}}

\subsection{Optimal Transport Construction}

We provide the detailed optimal-transport construction for Regime~B (large
simulation error, small domain gap).

Let $\cD_{\text{sim}}$ and $\cD_{\text{real}}$ be probability measures on
$\cS$ with finite second moments.  By the Brenier theorem~\cite{brenier1991},
there exists a convex potential $\varphi: \cS \to \R$ such that the optimal
transport map $T = \nabla\varphi$ satisfies $T_\# \cD_{\text{sim}} =
\cD_{\text{real}}$ and:
\begin{equation}
  W_2^2(\cD_{\text{sim}}, \cD_{\text{real}}) =
  \E_{\mathbf{s} \sim \cD_{\text{sim}}}\bigl[ \norm{T(\mathbf{s}) - \mathbf{s}}^2 \bigr].
\end{equation}

Now construct $M'(\mathbf{s}) = \Phi^*(\mathbf{s}) + \xi\mathbf{w}(\mathbf{s})$
where:
\begin{equation}
  \mathbf{w}(\mathbf{s}) =
  \frac{T(\Phi^*(\mathbf{s})) - \Phi^*(\mathbf{s})}{
    \E[\norm{T(\Phi^*(\mathbf{s})) - \Phi^*(\mathbf{s})}^2]^{1/2}}.
\end{equation}

This ensures that the perturbation has unit expected magnitude and is aligned
with the optimal transport direction.  The induced distribution of $M'$ is
$(M')_\# \cD_{\text{sim}} = \cD_{\text{real}}$ by construction, so
$\mathcal{E}_{\text{gap}} = 0$.

Meanwhile, the per-step simulation error is:
\begin{equation}
  \mathcal{E}_{\text{sim}}(M') =
  \E\bigl[ \norm{\xi \mathbf{w}(\mathbf{s})}^2 \bigr] = \xi^2.
\end{equation}

Setting $\xi^2 = \mathcal{E}_{\text{obs}}$ yields the claimed decomposition
with $\mathcal{E}_{\text{sim}} = \mathcal{E}_{\text{obs}}$ and
$\mathcal{E}_{\text{gap}} = 0$.

\subsection{Continuum of Solutions}

The non-uniqueness follows from the degrees of freedom in both constructions:
\begin{itemize}[nosep]
  \item Regime~A: The perturbation field $\eta(\mathbf{s})$ can be any
    divergence-free vector field on $\cS \setminus A$, yielding an
    infinite-dimensional family;
  \item Regime~B: The transport map $T$ is unique only for the quadratic cost,
    but alternative costs (e.g., asymmetric, regularized) yield distinct
    decompositions with the same $\mathcal{E}_{\text{obs}}$.
\end{itemize}

This establishes that the unidentifiability is not an artifact of a specific
construction but a fundamental property of the observational setting.
$\square$

%==============================================================================
%  ACKNOWLEDGMENTS
%==============================================================================

\section*{Acknowledgments}
We thank the NVIDIA Isaac Sim and Cosmos teams for making their platforms
available for research.  This work benefited from discussions with the
sim-to-real transfer community (\textbf{}) and the formal
verification community.  The Cercis naming honors the Eastern Redbud
(\emph{Cercis canadensis}), a tree known for its resilience across diverse
environments---a fitting metaphor for robust simulation auditing.

%==============================================================================
%  BIBLIOGRAPHY
%==============================================================================

\begin{thebibliography}{99}

\bibitem{cosmos2025}
NVIDIA Corporation.
\emph{NVIDIA Cosmos: A World Foundation Model Platform for Physical AI}.
Technical Report, 2025.

\bibitem{isaacsim2023}
NVIDIA Corporation.
\emph{Isaac Sim: Robotics Simulation and Synthetic Data Generation}.
NVIDIA Developer Documentation, 2023.

\bibitem{physx}
NVIDIA Corporation.
\emph{PhysX SDK 5.0 Documentation}.
NVIDIA GameWorks, 2023.

\bibitem{genie2024}
J.~Bruce, M.~Dennis, A.~Edwards, \emph{et al.}
\emph{Genie: Generative Interactive Environments}.
In \emph{ICML}, 2024.

\bibitem{unisim2023}
Y.~Yang, D.~Chen, S.~Tulsiani, \emph{et al.}
\emph{UniSim: A Neural Closed-Loop Sensor Simulator}.
In \emph{CVPR}, 2023.

\bibitem{tobin2017domain}
J.~Tobin, R.~Fong, A.~Ray, \emph{et al.}
\emph{Domain Randomization for Transferring Deep Neural Networks from
Simulation to the Real World}.
In \emph{IROS}, 2017.

\bibitem{peng2018sim}
X.~B.~Peng, M.~Andrychowicz, W.~Zaremba, and P.~Abbeel.
\emph{Sim-to-Real Transfer of Robotic Control with Dynamics Randomization}.
In \emph{ICRA}, 2018.

\bibitem{song2021score}
Y.~Song, J.~Sohl-Dickstein, D.~P.~Kingma, \emph{et al.}
\emph{Score-Based Generative Modeling through Stochastic Differential
Equations}.
In \emph{ICLR}, 2021.

\bibitem{villani2009}
C.~Villani.
\emph{Optimal Transport: Old and New}.
Springer, 2009.

\bibitem{basseville1993detection}
M.~Basseville and I.~V.~Nikiforov.
\emph{Detection of Abrupt Changes: Theory and Application}.
Prentice Hall, 1993.

\bibitem{page1954continuous}
E.~S.~Page.
\emph{Continuous Inspection Schemes}.
Biometrika, 41(1/2):100--115, 1954.

\bibitem{wald1945sequential}
A.~Wald.
\emph{Sequential Tests of Statistical Hypotheses}.
Annals of Mathematical Statistics, 16(2):117--186, 1945.

\bibitem{schulman2017ppo}
J.~Schulman, F.~Wolski, P.~Dhariwal, \emph{et al.}
\emph{Proximal Policy Optimization Algorithms}.
arXiv:1707.06347, 2017.

\bibitem{miyato2018spectral}
T.~Miyato, T.~Kataoka, M.~Koyama, and Y.~Yoshida.
\emph{Spectral Normalization for Generative Adversarial Networks}.
In \emph{ICLR}, 2018.

\bibitem{trockman2025mim}
A.~Trockman and J.~Z.~Kolter.
\emph{Mimetic Initialization for Lipschitz-Constrained Networks}.
In \emph{ICLR}, 2025.

\bibitem{pearl2009causality}
J.~Pearl.
\emph{Causality: Models, Reasoning, and Inference}.
Cambridge University Press, 2nd edition, 2009.

\bibitem{oberkampf2010verification}
W.~L.~Oberkampf and C.~J.~Roy.
\emph{Verification and Validation in Scientific Computing}.
Cambridge University Press, 2010.

\bibitem{hu2024eval}
Y.~Hu, \emph{et al.}
\emph{Evaluating World Models with Physical Benchmarks}.
In \emph{NeurIPS}, 2024.

\bibitem{michel2025genie}
M.~Michel, \emph{et al.}
\emph{Genie 2: Large-Scale Foundation World Models}.
DeepMind Technical Report, 2025.

\bibitem{hendrycks2019baseline}
D.~Hendrycks and K.~Gimpel.
\emph{A Baseline for Detecting Misclassified and Out-of-Distribution
Examples in Neural Networks}.
In \emph{ICLR}, 2019.

\bibitem{yang2022ood}
J.~Yang, K.~Zhou, Y.~Li, and Z.~Liu.
\emph{Generalized Out-of-Distribution Detection: A Survey}.
arXiv:2110.11334, 2022.

\bibitem{zhao2020sim}
W.~Zhao, J.~P.~Queralta, and T.~Westerlund.
\emph{Sim-to-Real Transfer in Deep Reinforcement Learning for Robotics:
A Survey}.
In \emph{SSCI}, 2020.

\bibitem{zhang2018crown}
H.~Zhang, T.-W.~Weng, P.-Y.~Chen, \emph{et al.}
\emph{Efficient Neural Network Robustness Certification with General
Activation Functions}.
In \emph{NeurIPS}, 2018.

\bibitem{wang2021beta}
S.~Wang, H.~Zhang, K.~Xu, \emph{et al.}
\emph{Beta-CROWN: Efficient Bound Propagation with Per-Neuron Split
Constraints for Neural Network Verification}.
In \emph{NeurIPS}, 2021.

\bibitem{spirtes2000causation}
P.~Spirtes, C.~Glymour, and R.~Scheines.
\emph{Causation, Prediction, and Search}.
MIT Press, 2nd edition, 2000.

\bibitem{brenier1991}
Y.~Brenier.
\emph{Polar Factorization and Monotone Rearrangement of Vector-Valued
Functions}.
Communications on Pure and Applied Mathematics, 44(4):375--417, 1991.

\bibitem{chang2019neural}
Y.-C.~Chang, N.~Roohi, and S.~Gao.
\emph{Neural Lyapunov Control}.
In \emph{NeurIPS}, 2019.

\bibitem{qin2021safe}
Z.~Qin, K.~Zhang, Y.~Chen, \emph{et al.}
\emph{Learning Safe Multi-Agent Control with Decentralized Neural Barrier
Certificates}.
In \emph{ICLR}, 2021.

\bibitem{waymo2023}
Waymo LLC.
\emph{Waymo Safety Report: 2023 Update}.
Technical Report, 2023.

\bibitem{fazlyab2019efficient}
M.~Fazlyab, A.~Roby, V.~M.~Preciado, and G.~J.~Pappas.
\emph{Efficient and Accurate Estimation of Lipschitz Constants for Deep
Neural Networks}.
In \emph{NeurIPS}, 2019.

\bibitem{katz2017reluplex}
G.~Katz, C.~Barrett, D.~L.~Dill, \emph{et al.}
\emph{Reluplex: An Efficient SMT Solver for Verifying Deep Neural Networks}.
In \emph{CAV}, 2017.

\bibitem{lamport1982byzantine}
L.~Lamport, R.~Shostak, and M.~Pease.
\emph{The Byzantine Generals Problem}.
ACM Transactions on Programming Languages and Systems, 4(3):382--401, 1982.

\bibitem{blanchard2017machine}
P.~Blanchard, E.~M.~El Mhamdi, R.~Guerraoui, and J.~Stainer.
\emph{Machine Learning with Adversaries: Byzantine Tolerant Gradient
Descent}.
In \emph{NeurIPS}, 2017.

\end{thebibliography}

\end{document}
